% ============================================================
% 04-bezier-matrix-straight-code-explained.tex
% Compile with:  pdflatex -shell-escape 04-bezier-matrix-straight-code-explained.tex
% Run twice for table of contents.
% ============================================================

\documentclass[12pt, a4paper]{article}

% ---- Encoding & fonts ----------------------------------------
\usepackage[T1]{fontenc}
\usepackage[utf8]{inputenc}
\usepackage{lmodern}

% ---- Page geometry -------------------------------------------
\usepackage[
  top=2.5cm, bottom=2.5cm,
  left=2.8cm, right=2.8cm
]{geometry}

% ---- Colours -------------------------------------------------
\usepackage[dvipsnames]{xcolor}
\definecolor{codebg}{RGB}{248,248,248}
\definecolor{rulecolor}{RGB}{180,180,180}
\definecolor{examplebg}{RGB}{240,248,255}
\definecolor{examplerule}{RGB}{100,149,237}
\definecolor{notebg}{RGB}{255,250,230}
\definecolor{noterule}{RGB}{210,180,50}
\definecolor{samebg}{RGB}{245,245,245}
\definecolor{samerule}{RGB}{160,160,160}

% ---- Minted --------------------------------------------------
\usepackage{minted}
\setminted[julia]{
  bgcolor    = codebg,
  frame      = leftline,
  framerule  = 1.5pt,
  rulecolor  = rulecolor,
  linenos    = false,
  fontsize   = \small,
  breaklines = true,
  tabsize    = 4,
  style      = friendly,
}
\setminted[text]{
  bgcolor    = codebg,
  fontsize   = \small,
  breaklines = true,
}

% ---- Math ----------------------------------------------------
\usepackage{amsmath}
\usepackage{amssymb}

% ---- TikZ ----------------------------------------------------
\usepackage{tikz}
\usetikzlibrary{shapes.geometric, arrows.meta, positioning, backgrounds}

% ---- Boxes ---------------------------------------------------
\usepackage{mdframed}

\newmdenv[
  backgroundcolor  = examplebg,
  linecolor        = examplerule, linewidth=1.5pt,
  leftline=true, rightline=false, topline=false, bottomline=false,
  innerleftmargin=10pt, innerrightmargin=8pt,
  innertopmargin=6pt, innerbottommargin=6pt,
  skipabove=6pt, skipbelow=4pt,
]{examplebox}

\newmdenv[
  backgroundcolor  = notebg,
  linecolor        = noterule, linewidth=1.5pt,
  leftline=true, rightline=false, topline=false, bottomline=false,
  innerleftmargin=10pt, innerrightmargin=8pt,
  innertopmargin=6pt, innerbottommargin=6pt,
  skipabove=6pt, skipbelow=4pt,
]{notebox}

\newmdenv[
  backgroundcolor  = samebg,
  linecolor        = samerule, linewidth=1.5pt,
  leftline=true, rightline=false, topline=false, bottomline=false,
  innerleftmargin=10pt, innerrightmargin=8pt,
  innertopmargin=6pt, innerbottommargin=6pt,
  skipabove=6pt, skipbelow=4pt,
]{samebox}

% ---- Hyperlinks ----------------------------------------------
\usepackage[colorlinks=true, linkcolor=black, urlcolor=MidnightBlue]{hyperref}

% ---- Misc ----------------------------------------------------
\usepackage{booktabs}
\usepackage{needspace}
\usepackage{parskip}
\usepackage{titlesec}
\usepackage{fancyhdr}
\usepackage{datetime2}

% ---- Section formatting --------------------------------------
\titleformat{\section}
  {\large\bfseries\color{MidnightBlue}}
  {\thesection.}{0.6em}{}[\vspace{-4pt}\rule{\linewidth}{0.4pt}]
\titleformat{\subsection}
  {\normalsize\bfseries\color{BrickRed}}
  {\thesubsection.}{0.5em}{}
\titleformat{\subsubsection}
  {\normalsize\bfseries}
  {\thesubsubsection.}{0.5em}{}

% ---- Header / footer -----------------------------------------
\pagestyle{fancy}
\fancyhf{}
\rhead{\textcolor{gray}{\small 04-bezier-matrix-straight-code-explained}}
\lhead{\textcolor{gray}{\small B\'ezier Matrix Straight --- Code Explanation}}
\cfoot{\thepage}
\renewcommand{\headrulewidth}{0.3pt}

% ---- Helpers -------------------------------------------------
\newcommand{\cd}[1]{\texttt{#1}}
\newcommand{\SAME}{\textcolor{samerule}{\textbf{[Same as script~03]}}\;}

% ---- Title ---------------------------------------------------
\title{%
  \vspace{1cm}
  {\LARGE\bfseries B\'ezier Curve --- Matrix Method (Straight)}\\[0.4em]
  {\large $P(u) = U \cdot M \cdot G$}\\[0.4em]
  {\large Code Explanation for Julia Beginners}\\[0.4em]
  {\normalsize\texttt{04-bezier-matrix-straight.jl}}
}
\author{E.R.\ Nurwijayadi}
\date{\DTMtoday}


% =============================================================
\begin{document}
% =============================================================

\maketitle
\thispagestyle{empty}

\begin{center}
\textit{Directed and curated by E.R.\ Nurwijayadi.
Code and explanations AI-generated.}
\end{center}

\begin{center}
\begin{minipage}{0.85\textwidth}
\small
This document explains the Julia code in
\cd{04-bezier-matrix-straight.jl} for readers new to Julia.
The majority of this script is identical to script~03.
This document focuses on what changed: the \cd{section()}
helper that replaces \cd{print\_step\_header()}, the five
new focused print functions, and one small change to
\cd{print\_labeled\_matrix}.
Shared components are listed with a reference to
\cd{03-bezier-matrix-pedagogical-code-explained}
rather than repeated.
\end{minipage}
\end{center}

\vspace{1em}\hrule\vspace{1.5em}
\tableofcontents
\newpage


% =============================================================
\needspace{8\baselineskip}
\section{Type Map: How the Pieces Fit Together}
% =============================================================

\textit{The computation core is unchanged from script~03.
Only the print layer changed --- the five numbered step
functions were replaced by five focused, single-purpose
matrix printers.}

\bigskip

\tikzset{
  structbox/.style={
    rectangle, rounded corners=4pt,
    draw=MidnightBlue, line width=1.5pt,
    fill=examplebg,
    text width=3.4cm, align=center,
    font=\small\bfseries, inner sep=8pt
  },
  funcbox/.style={
    rectangle, rounded corners=2pt,
    draw=gray!60, line width=1pt,
    fill=white,
    text width=3.4cm, align=center,
    font=\small, inner sep=7pt
  },
  samefunc/.style={
    rectangle, rounded corners=2pt,
    draw=samerule, line width=1pt,
    fill=samebg,
    text width=3.4cm, align=center,
    font=\small, inner sep=7pt
  },
  entrybox/.style={
    rectangle, rounded corners=4pt,
    draw=BrickRed, line width=1.5pt,
    fill=notebg,
    text width=2.4cm, align=center,
    font=\small\bfseries, inner sep=8pt
  },
  myarrow/.style={-Stealth, thick, draw=gray!70},
  bluearrow/.style={-Stealth, thick, draw=MidnightBlue},
  redarrow/.style={-Stealth, thick, draw=BrickRed!80},
  arrowlabel/.style={font=\footnotesize\itshape, fill=white, inner sep=2pt}
}

\begin{center}
\begin{tikzpicture}[node distance=1.4cm and 2.8cm]

  % -- BezierMatrix struct (unchanged) ----------------------
  \node[structbox] (bm) {%
    \texttt{BezierMatrix}\\[3pt]
    \footnotesize\normalfont
    \texttt{G, M, MG::Matrix\{Float64\}}\\
    \texttt{axes, n, degree, dim}
  };

  % -- derive_M + evaluate_all (unchanged) ------------------
  \node[samefunc, below left=1.6cm and 0.3cm of bm] (core) {%
    \texttt{derive\_M()}\\
    \texttt{evaluate\_all()}\\[2pt]
    \footnotesize\normalfont
    \textit{unchanged from script~03}
  };

  % -- fmt helpers (unchanged) ------------------------------
  \node[samefunc, below right=1.6cm and 0.3cm of bm] (fmt) {%
    \textbf{Formatting helpers}\\[2pt]
    \footnotesize\normalfont
    \texttt{fmt}, \texttt{fmt\_frac},
    \texttt{fmt\_poly}\\
    \textit{unchanged from script~03}
  };

  % -- table printers (mostly unchanged) --------------------
  \node[funcbox, below=1.3cm of core] (tbl) {%
    \textbf{Table printers}\\[2pt]
    \footnotesize\normalfont
    \texttt{print\_table} (unchanged)\\
    \texttt{print\_labeled\_matrix}\\
    \textit{(lost \texttt{style} kwarg)}
  };

  % -- NEW: section() helper --------------------------------
  \node[funcbox, below=1.3cm of fmt] (sec) {%
    \texttt{section()} \textbf{[new]}\\[2pt]
    \footnotesize\normalfont
    \textit{replaces}\\
    \texttt{print\_step\_header()}
  };

  % -- NEW: focused print functions -------------------------
  \node[funcbox, below=1.3cm of tbl, xshift=2.1cm] (print) {%
    \textbf{Print functions [new]}\\[2pt]
    \footnotesize\normalfont
    \texttt{print\_G}, \texttt{print\_M}\\
    \texttt{print\_MG}, \texttt{print\_polynomials}\\
    \texttt{print\_eval\_table}\\
    \texttt{print\_all}
  };

  % -- main (unchanged) -------------------------------------
  \node[entrybox, right=2.2cm of bm] (main) {%
    \texttt{main()}\\[2pt]
    \footnotesize\normalfont\normalcolor
    \textit{entry point}\\
    \textit{unchanged}
  };

  % -- Arrows -----------------------------------------------
  \draw[bluearrow] (core)  -- node[arrowlabel,left]{builds $M$, $MG$} (bm);
  \draw[myarrow]   (bm)    -- node[arrowlabel,above left]{calls} (core);
  \draw[myarrow]   (print) -- node[arrowlabel,left]{calls}  (tbl);
  \draw[myarrow]   (print) -- node[arrowlabel,right]{calls} (sec);
  \draw[myarrow]   (print) -- node[arrowlabel,above]{reads} (bm);
  \draw[myarrow]   (print) |- node[arrowlabel,right,pos=0.8]{calls} (core);
  \draw[myarrow]   (print) -- node[arrowlabel,above]{calls} (fmt);
  \draw[redarrow]  (main)  -- node[arrowlabel,above]{calls} (bm);
  \draw[redarrow]  (main)  |- node[arrowlabel,right,pos=0.75]{calls} (print);

\end{tikzpicture}
\end{center}

\bigskip
\begin{notebox}
\textbf{Legend:}
Grey-shaded boxes are components unchanged from script~03.
White boxes are new or modified.
The computation core (\cd{BezierMatrix}, \cd{derive\_M},
\cd{evaluate\_all}, formatting helpers, \cd{print\_table})
is entirely reused.
Only the print layer changed.
\end{notebox}


% =============================================================
\needspace{8\baselineskip}
\section{What Is Identical to Script 03}
% =============================================================

The following components are copied verbatim from script~03.
They are documented fully in
\cd{03-bezier-matrix-pedagogical-code-explained}.

\begin{samebox}
\begin{tabular}{lp{9cm}}
\toprule
Component & Notes \\
\midrule
\cd{CASE\_DATA} & Same control points, same $u$ values. \\
\cd{BezierMatrix} struct & Fields \cd{G}, \cd{M}, \cd{MG},
  \cd{axes}, \cd{n}, \cd{degree}, \cd{dim} --- unchanged. \\
\cd{BezierMatrix} constructor & Calls \cd{derive\_M},
  computes \cd{M * points}. \\
\cd{derive\_M} & Double loop, index arithmetic
  \cd{M[n-k+1, i+1]}, \cd{zeros()} initialisation. \\
\cd{evaluate\_all} & 2-D array comprehension,
  one batch matrix multiply \cd{Uall * bm.MG}. \\
\cd{fmt} & Integer or 4 d.p. \\
\cd{fmt\_frac} & Float $\to$ rational string,
  short-circuit \cd{||} guard. \\
\cd{fmt\_poly} & Coefficient vector $\to$ polynomial string,
  reversed power order. \\
\cd{print\_table} & Keyword args, \cd{Union\{...,Nothing\}},
  \cd{:simple}/\cd{:outline} styles. \\
\cd{ruler} & 66-character divider. \\
\cd{SUPERSCRIPTS} & Superscript character dict. \\
\cd{main} & \cd{CASE\_DATA[CASE]} $\to$ \cd{BezierMatrix}
  $\to$ \cd{print\_all}. \\
\bottomrule
\end{tabular}
\end{samebox}


% =============================================================
\needspace{8\baselineskip}
\section{What Changed: \texttt{print\_labeled\_matrix}}
% =============================================================

Script~03:
\begin{minted}{julia}
function print_labeled_matrix(
        mat        ::Matrix{Float64},
        row_labels ::Vector{String},
        col_labels ::Vector{String};
        fmt_fn     = fmt,
        style      ::Symbol = :simple)
    rows = [[fmt_fn(mat[r,c]) for c in 1:size(mat,2)]
            for r in 1:size(mat,1)]
    print_table(rows, col_labels; row_labels=row_labels, style=style)
end
\end{minted}

Script~04:
\begin{minted}{julia}
function print_labeled_matrix(
        mat        ::Matrix{Float64},
        row_labels ::Vector{String},
        col_labels ::Vector{String};
        fmt_fn     = fmt)
    rows = [[fmt_fn(mat[r,c]) for c in 1:size(mat,2)]
            for r in 1:size(mat,1)]
    print_table(rows, col_labels; row_labels=row_labels)
end
\end{minted}

The \cd{style} keyword argument is gone.
Script~03 needed it because \cd{print\_labeled\_matrix}
was called from both simple-table contexts (matrix display)
and potentially outline contexts.
Script~04 only ever calls \cd{print\_labeled\_matrix} for
the three matrices ($G$, $M$, $M \cdot G$), which all use
the default \cd{:simple} style.
The outline table is only used in \cd{print\_eval\_table},
which calls \cd{print\_table} directly with
\cd{style=:outline}.

Removing an unused parameter is good practice --- it
simplifies the interface and removes a source of confusion.


% =============================================================
\needspace{8\baselineskip}
\section{What Is New: \texttt{section()}}
% =============================================================

\begin{minted}{julia}
section(title::String) = (
    println();
    println("  $title");
    println("  " * "-"^(WIDTH-2))
)
\end{minted}

This replaces \cd{print\_step\_header(step::Int, title::String)}
from script~03.

The difference is small but meaningful:

\begin{center}
\begin{tabular}{lll}
\toprule
 & Script~03 & Script~04 \\
\midrule
Signature & \cd{(step::Int, title::String)} &
  \cd{(title::String)} \\
Output & \cd{STEP 2 --- Derive Basis Matrix...} &
  \cd{M --- Basis Matrix (derived from Bernstein)} \\
\end{tabular}
\end{center}

Script~03 numbered its steps because the output was a
structured walkthrough.
Script~04 has no steps --- each section is an independent
block labelled by its content, not by a sequence number.

\begin{examplebox}
\textbf{Script~03:}
\begin{minted}{text}
  STEP 2 - Derive Basis Matrix M from Bernstein Polynomials
  ----------------------------------------------------------------
\end{minted}

\textbf{Script~04:}
\begin{minted}{text}
  M  - Basis Matrix (derived from Bernstein)
  ----------------------------------------------------------------
\end{minted}

Same dashed rule, same indentation.
Just no step number, and the title directly names the object
being displayed.
\end{examplebox}


% =============================================================
\needspace{8\baselineskip}
\section{What Is New: The Five Print Functions}
% =============================================================

\textit{Script~03 had five numbered \cd{print\_stepN} functions,
each performing a major computation or derivation and
printing detailed working.
Script~04 replaces them with five focused functions ---
each one prints exactly one thing.}

% -------------------------------------------------------
\needspace{6\baselineskip}
\subsection{\texttt{print\_G}}

\begin{minted}{julia}
function print_G(bm::BezierMatrix)
    section("G  - Geometry Matrix (Control Points)")
    println()
    print_labeled_matrix(bm.G, ["P$(i-1)" for i in 1:bm.n], bm.axes)
end
\end{minted}

\textbf{Input:} reads \cd{bm.G}, \cd{bm.n}, \cd{bm.axes}.

\textbf{What it does:} prints a section header, then calls
\cd{print\_labeled\_matrix} with row labels
\cd{["P0","P1","P2","P3","P4"]} and column labels from
\cd{bm.axes}.

\textbf{Output:}
\begin{minted}{text}
  G  - Geometry Matrix (Control Points)
  ----------------------------------------------------------------

      x  y  z
      -  -  -
  P0  0  0  1
  P1  0  4  1
  P2  4  0  1
  P3  4  4  1
  P4  5  4  1
\end{minted}

% -------------------------------------------------------
\needspace{6\baselineskip}
\subsection{\texttt{print\_M}}

\begin{minted}{julia}
function print_M(bm::BezierMatrix)
    section("M  - Basis Matrix (derived from Bernstein)")
    n = bm.degree
    println()
    print_labeled_matrix(bm.M,
        ["u^$(n-k)" for k in 0:n],
        ["P$(i-1)" for i in 1:bm.n])
end
\end{minted}

\textbf{Input:} reads \cd{bm.M}, \cd{bm.degree}, \cd{bm.n}.

\textbf{What it does:} prints the basis matrix with row
labels \cd{["u\^{}4","u\^{}3","u\^{}2","u\^{}1","u\^{}0"]}
(built by the comprehension \cd{["u\^{}\$(n-k)" for k in 0:n]})
and column labels \cd{["P0",..."P4"]}.

The column labels here are the control point names,
not the axis names --- because columns of $M$ correspond
to Bernstein basis polynomials indexed by control point,
not to spatial axes.

\textbf{Output:}
\begin{minted}{text}
  M  - Basis Matrix (derived from Bernstein)
  ----------------------------------------------------------------

       P0   P1   P2  P3  P4
       --  ---  ---  --  --
  u^4   1   -4    6  -4   1
  u^3  -4   12  -12   4   0
  u^2   6  -12    6   0   0
  u^1  -4    4    0   0   0
  u^0   1    0    0   0   0
\end{minted}

% -------------------------------------------------------
\needspace{6\baselineskip}
\subsection{\texttt{print\_MG}}

\begin{minted}{julia}
function print_MG(bm::BezierMatrix)
    section("M * G  - Polynomial Coefficient Matrix")
    n = bm.degree
    println()
    print_labeled_matrix(bm.MG,
        ["u^$(n-k)" for k in 0:n],
        bm.axes)
end
\end{minted}

\textbf{Input:} reads \cd{bm.MG}, \cd{bm.degree}, \cd{bm.axes}.

\textbf{What it does:} prints the pre-computed $M \cdot G$
matrix with row labels for powers of $u$ and column labels
from \cd{bm.axes} (the spatial axes $x$, $y$, $z$).

Note the column labels here switch back to \cd{bm.axes},
because columns of $M \cdot G$ correspond to axes,
not to control points.

\textbf{Output:}
\begin{minted}{text}
  M * G  - Polynomial Coefficient Matrix
  ----------------------------------------------------------------

         x    y  z
       ---  ---  -
  u^4   13  -28  0
  u^3  -32   64  0
  u^2   24  -48  0
  u^1    0   16  0
  u^0    0    0  1
\end{minted}

% -------------------------------------------------------
\needspace{6\baselineskip}
\subsection{\texttt{print\_polynomials}}

\begin{minted}{julia}
function print_polynomials(bm::BezierMatrix)
    section("Resulting Polynomials")
    println()
    for (d, axis) in enumerate(bm.axes)
        println("  $(axis)(u)  =  $(fmt_poly(bm.MG[:,d], bm.degree))")
    end
end
\end{minted}

\textbf{Input:} reads \cd{bm.MG}, \cd{bm.degree}, \cd{bm.axes}.

\textbf{What it does:} loops over axes.
For each axis \cd{d}, takes the column slice \cd{bm.MG[:,d]}
--- a \cd{Vector\{Float64\}} of monomial coefficients from
highest power down --- and passes it to \cd{fmt\_poly}
along with \cd{bm.degree} to produce the polynomial string.

\begin{examplebox}
\textbf{For axis \cd{"x"} (\cd{d=1}):}

\cd{bm.MG[:,1] = [13.0, -32.0, 24.0, 0.0, 0.0]}

\cd{fmt\_poly([13.0, -32.0, 24.0, 0.0, 0.0], 4)}
$\to$ \cd{"13u\^{}4  - 32u\^{}3  + 24u\^{}2"}

Printed as: \cd{x(u)  =  13u\^{}4  - 32u\^{}3  + 24u\^{}2}
\end{examplebox}

\textbf{Output:}
\begin{minted}{text}
  Resulting Polynomials
  ----------------------------------------------------------------

  x(u)  =  13u^4  - 32u^3  + 24u^2
  y(u)  =  -28u^4  + 64u^3  - 48u^2  + 16u
  z(u)  =  1
\end{minted}

% -------------------------------------------------------
\needspace{6\baselineskip}
\subsection{\texttt{print\_eval\_table}}

\begin{minted}{julia}
function print_eval_table(bm::BezierMatrix,
                          u_values::Vector{Float64})
    section("Evaluation Table  -  P(u) = U * (M * G)")
    pts  = evaluate_all(bm, u_values)
    rows = Vector{Vector{String}}()
    for (i, u) in enumerate(u_values)
        row = [fmt(u)]
        for d in 1:bm.dim
            v = pts[i, d]
            push!(row, "$(fmt_frac(v))  ($(@sprintf("%.4f", v)))")
        end
        push!(rows, row)
    end
    println()
    print_table(rows,
        vcat(["u"], ["$(ax)(u)" for ax in bm.axes]);
        style=:outline)
end
\end{minted}

\textbf{Input:} reads \cd{bm.dim}, \cd{bm.axes}, and
the external \cd{u\_values} vector.

\textbf{What it does:} identical logic to
\cd{print\_step5\_table} in script~03.
Calls \cd{evaluate\_all} for the batch result, builds
the rows vector with \cd{fmt\_frac} + \cd{@sprintf} cells,
and prints with \cd{style=:outline}.

The only difference from script~03 is the section header:
\cd{section("Evaluation Table ...")} instead of
\cd{print\_step\_header(5, "Summary ...")}.

\textbf{Output:}
\begin{minted}{text}
  Evaluation Table  -  P(u) = U * (M * G)
  ----------------------------------------------------------------

  +=============================================================+
  | u       x(u)               y(u)               z(u)        |
  |=============================================================|
  | 0       0  (0.0000)        0  (0.0000)   1  (1.0000)      |
  | 0.2500  269//256  (1.0508) 121//64 (1.8906) 1  (1.0000)   |
  | 0.5000  45//16  (2.8125)   9//4  (2.2500)   1  (1.0000)   |
  | 0.7500  1053//256 (4.1133) 201//64 (3.1406) 1  (1.0000)   |
  | 1       5  (5.0000)        4  (4.0000)       1  (1.0000)  |
  +=============================================================+
\end{minted}

% -------------------------------------------------------
\needspace{6\baselineskip}
\subsection{\texttt{print\_all}}

\begin{minted}{julia}
function print_all(bm::BezierMatrix,
                   u_values::Vector{Float64}, case_num::Int)
    print_title(bm, case_num)
    print_G(bm)
    print_M(bm)
    print_MG(bm)
    print_polynomials(bm)
    print_eval_table(bm, u_values)
    println()
    println(ruler())
    println()
end
\end{minted}

\textbf{What it does:} orchestrates the five print functions
in order, with title and closing ruler.

Compare to script~03's \cd{print\_all}:

\begin{center}
\begin{tabular}{ll}
\toprule
Script~03 & Script~04 \\
\midrule
\cd{print\_title} & \cd{print\_title} \\
\cd{print\_control\_points} & \textit{(replaced by \cd{print\_G})} \\
\cd{print\_step1\_define} & \textit{(dropped --- no formulation block)} \\
\cd{print\_step2\_derive\_M} & \cd{print\_M} \\
\cd{print\_step3\_MG} & \cd{print\_MG} + \cd{print\_polynomials} \\
\cd{print\_step4\_evaluate} & \textit{(dropped --- no dot product display)} \\
\cd{print\_step5\_table} & \cd{print\_eval\_table} \\
\bottomrule
\end{tabular}
\end{center}

\begin{notebox}
\textbf{Why split \cd{print\_step3\_MG} into two functions?}

In script~03, \cd{print\_step3\_MG} printed both the
$M \cdot G$ matrix and the resulting polynomial strings
in one block under the same step header.
Script~04 separates them: \cd{print\_MG} prints the matrix,
\cd{print\_polynomials} prints the strings, each with its
own section header.
This gives the polynomials their own visual prominence ---
they are the practical result the user cares about most.
\end{notebox}


% =============================================================
\needspace{8\baselineskip}
\section{Design Principle: One Function, One Thing}
% =============================================================

\textit{The real lesson of this script is not in any single
function --- it is in how the functions are organised.}

Script~03 organised its output around the \emph{process}:
Step~1 defines, Step~2 derives, Step~3 computes, and so on.
Each step function did multiple things: explained, computed,
and printed.

Script~04 organises its output around the \emph{data}:
one function per matrix, one function for the polynomials,
one function for the table.
Each print function does exactly one thing: print one object.

This is a principle called \textbf{separation of concerns}.
Each function has a single, clear responsibility.
If you want to change how the $M$ matrix is displayed,
you edit \cd{print\_M} and nothing else.
If you want to add a new section, you write a new function
and add one line to \cd{print\_all}.

\begin{examplebox}
\textbf{Adding a new section would be trivial.}
Suppose you wanted to print the $U$ vector for $u = 0.5$
as a new section. You would write:

\begin{minted}{julia}
function print_U(bm::BezierMatrix, u::Float64)
    section("U  - Parameter Vector at u = $u")
    n = bm.degree
    Uv = [u^(n-k) for k in 0:n]
    println()
    # ... format and print Uv
end
\end{minted}

Then add \cd{print\_U(bm, 0.5)} to \cd{print\_all}.
No other function needs to change.
\end{examplebox}

\vspace{2em}\hrule\vspace{0.5em}
{\small\textcolor{gray}{%
Code explanation for \texttt{04-bezier-matrix-straight.jl}
\textbar{} Directed and curated by E.R.\ Nurwijayadi
\textbar{} Code and explanations AI-generated}}

\end{document}
